\chapter{pytop'a Hızlı Giriş}

pytop kurulumu, temel uzay nesneleri, \texttt{Result} tipi ve tag sistemi —
kılavuzun geri kalanını okumadan önce bu bölümü çalıştırın.

%-------------------------------------------------------------------
\section{Kurulum ve Import}

\begin{lstlisting}[language=Python]
pip install -e .   # git kokundan
\end{lstlisting}

Her bolumde ihtiyac duyulan semboller dogrudan \texttt{pytop}'tan import edilir.
Bu bolum bir \emph{tur} oldugundan, ilk hucrede sonraki tum hucrelerin
kullanacagi sembolleri tek seferde ice aktariyoruz:

\begin{lstlisting}[language=Python]
import pytop
from pytop import (
    discrete_topology, indiscrete_topology, sierpinski_space, make_topology,
    is_compact, is_connected, is_t0, is_t1, is_t2,
    count_topologies_on_n_points,
    simplicial_complex, simplicial_homology, betti_numbers,
    FiniteMetricSpace, persistent_homology, persistence_betti_numbers,
)
from pytop.experimental.spaces import finite_circle, pi1_space
print("pytop surumu:", pytop.__version__)
\end{lstlisting}

\begin{verbatim}
pytop surumu: 1.7.0
\end{verbatim}

\begin{sezgi}
pytop'u iki katlı bir bina gibi düşünün. Alt kat \textbf{betimsel}
(descriptive) katmandır: ünlü uzayların ve teoremlerin kayıtlı, kaynaklı
bilgisini tutar --- bilir ama hesaplamaz. Üst kat \textbf{yapıcı}
(constructive) katmandır: ham girdiden invaryantları \emph{hesaplar} ---
homoloji, kalıcı homoloji, düğüm polinomları. Bu bölüm her iki kata da
kısa bir tur attırır.
\end{sezgi}

\begin{figure}[ht]
\centering
\begin{tikzpicture}[font=\small, node distance=6mm,
    kat/.style={draw, rounded corners, minimum width=78mm, minimum height=15mm,
                align=center, inner sep=5pt},
    kut/.style={draw, rounded corners, fill=white, inner sep=3pt,
                align=center, font=\scriptsize}]
  \node[kat, fill=green!8] (desc) {};
  \node[anchor=north west, font=\bfseries] at (desc.north west) {Betimsel katman};
  \node[anchor=south east, font=\scriptsize\itshape] at (desc.south east) {bilir, hesaplamaz};
  \node[kut, below=2mm of desc.north, yshift=-4mm] (d1) {\texttt{*Profile} kayitlari};
  \node[kut, right=4mm of d1] (d2) {\texttt{get\_*\_profiles()}};
  \node[kut, left=4mm of d1] (d0) {unlu uzay/teorem};

  \node[kat, fill=blue!8, below=14mm of desc] (cons) {};
  \node[anchor=north west, font=\bfseries] at (cons.north west) {Yapici katman};
  \node[anchor=south east, font=\scriptsize\itshape] at (cons.south east) {ham girdiden hesaplar};
  \node[kut, below=2mm of cons.north, yshift=-4mm] (c1) {\texttt{simplicial\_homology}};
  \node[kut, right=4mm of c1] (c2) {\texttt{persistent\_homology}};
  \node[kut, left=4mm of c1] (c0) {\texttt{jones\_polynomial}};

  \draw[-{Stealth[length=2.6mm]}, thick]
    (cons.north) -- (desc.south)
    node[midway, right=1mm, font=\scriptsize, align=left]
    {hesaplanan sonuc\\ kayitlari dogrular};
\end{tikzpicture}

\caption{pytop'un iki katmanı: betimsel kayıtlar ile yapıcı motorlar.}
\end{figure}

%-------------------------------------------------------------------
\section{Uzay Nesneleri}

\begin{table}[h]
\centering
\begin{tabular}{lll}
\toprule
\textbf{Kurucu} & \textbf{Topoloji} & \textbf{Sonuç} \\
\midrule
\texttt{discrete\_topology(*pts)} & Her altküme açık & Ayrık \\
\texttt{indiscrete\_topology(*pts)} & Yalnız $\emptyset$ ve $X$ açık & İndiscrete \\
\texttt{sierpinski\_space()} & $\{\emptyset, \{1\}, \{0,1\}\}$ & Sierpiński \\
\texttt{make\_topology(carrier, *open\_sets)} & Kullanıcı tanımlı & Özel \\
\bottomrule
\end{tabular}
\caption{Temel uzay kurucuları}
\end{table}

\begin{lstlisting}[language=Python]
d = discrete_topology(0, 1, 2)
print("carrier:", sorted(d.carrier))
print("topoloji boyutu:", len(d.topology))
print("etiketler:", sorted(d.tags))
\end{lstlisting}

\begin{verbatim}
carrier: [0, 1, 2]
topoloji boyutu: 8
etiketler: ['discrete', 'finite', 'hausdorff', 'metrizable', 'normal', 'regular']
\end{verbatim}

%-------------------------------------------------------------------
\section{Result Tipi}

pytop'un çoğu yüklemi bir \texttt{Result} nesnesi döner.

\begin{lstlisting}[language=Python]
s = sierpinski_space()
r = is_t2(s)
print(".status:", r.status)        # 'true' | 'false' | 'unknown'
print(".value:", r.value)
print(".mode:", r.mode)
print(".justification:", r.justification)
\end{lstlisting}

\begin{verbatim}
.status: false
.value: hausdorff
.mode: exact
.justification: ['The explicit finite topology fails hausdorff.']
\end{verbatim}

\texttt{.status} her zaman \texttt{'true'}, \texttt{'false'} veya \texttt{'unknown'}
dizesidir --- Python bool'u değil.

%-------------------------------------------------------------------
\section{Temel Uzaylar Turu}

\begin{lstlisting}[language=Python]
spaces = {
    "Ayrik D(0,1,2)":     discrete_topology(0, 1, 2),
    "Indiscrete I(0,1,2)": indiscrete_topology(0, 1, 2),
    "Sierpinski S":        sierpinski_space(),
}
for name, sp in spaces.items():
    print(name,
          is_compact(sp).status,
          is_connected(sp).status,
          is_t2(sp).status)
\end{lstlisting}

\begin{verbatim}
Ayrik D(0,1,2)      true  false  true
Indiscrete I(0,1,2) true  true   false
Sierpinski S        true  true   false
\end{verbatim}

%-------------------------------------------------------------------
\section{Özel Topoloji: make\_topology}

\begin{lstlisting}[language=Python]
X = [1, 2, 3, 4]
tau = [set(), {1,2}, {3,4}, {1,2,3,4}]
fts = make_topology(X, *tau)
print("topoloji boyutu:", len(fts.topology))
print("kompakt:", is_compact(fts).status)
print("T2:", is_t2(fts).status)
\end{lstlisting}

\begin{verbatim}
topoloji boyutu: 4
kompakt: true
T2: false
\end{verbatim}

%-------------------------------------------------------------------
\section{Tag Sistemi}

\begin{lstlisting}[language=Python]
print("Ayrik:    ", sorted(discrete_topology(0,1,2).tags))
print("Indiscrete:", sorted(indiscrete_topology(0,1,2).tags))
print("Sierpinski:", sorted(sierpinski_space().tags))
\end{lstlisting}

\begin{verbatim}
Ayrik:      ['discrete', 'finite', 'hausdorff', 'metrizable', 'normal', 'regular']
Indiscrete: ['compact', 'connected', 'finite', 'indiscrete']
Sierpinski: ['compact', 'connected', 'finite', 't0']
\end{verbatim}

Etiketler yüklem fonksiyonlarının karar mekanizmasını hızlandırır:
\texttt{is\_t2(d)} açık küme taraması yerine \texttt{'hausdorff' in d.tags}
kontrolüyle sonuç verir.

%-------------------------------------------------------------------
\section{n Nokta Üzerindeki Topoloji Sayısı}

\begin{lstlisting}[language=Python]
for n in range(1, 6):
    print(f"n={n}: {count_topologies_on_n_points(n)}")
\end{lstlisting}

\begin{verbatim}
n=1: 1
n=2: 4
n=3: 29
n=4: 355
n=5: 6942
\end{verbatim}

Topoloji sayısı hızla büyür; sonlu uzayların tüm topolojileri üzerinde kaba
kuvvet taraması yalnızca küçük $n$ için pratiktir.

%-------------------------------------------------------------------
\section{pytop Ne Hesaplar? — Genel Harita}

Önceki bölümler \textbf{betimsel/sonlu} katmanı gezdirdi. pytop'un asıl gücü
\textbf{yapıcı çekirdektir}: ham bir kompleks ya da nokta bulutundan cebirsel
invaryantları gerçekten \emph{hesaplar}.

\begin{figure}[ht]
\centering
\begin{tikzpicture}[font=\small, node distance=10mm,
    gir/.style={draw, rounded corners, fill=orange!12, align=center, inner sep=5pt},
    inv/.style={draw, rounded corners, fill=blue!8, align=center, inner sep=5pt}]
  \node[gir] (pts) {nokta bulutu};
  \node[gir, below=8mm of pts] (cpx) {simpleks kompleks};
  \node[gir, below=8mm of cpx] (knot) {dugum diyagrami};

  \node[inv, right=34mm of pts] (ph) {kalici homoloji\\ \scriptsize barkod, $H_*$};
  \node[inv, right=34mm of cpx] (hom) {homoloji\\ \scriptsize Betti, torsiyon};
  \node[inv, right=34mm of knot] (jp) {Jones / Alexander\\ \scriptsize polinom};

  \draw[-{Stealth[length=2.6mm]}, thick] (pts)  -- (ph)
        node[midway, above, font=\scriptsize] {Vietoris--Rips};
  \draw[-{Stealth[length=2.6mm]}, thick] (cpx)  -- (hom)
        node[midway, above, font=\scriptsize] {sinir matrisi $\to$ SNF};
  \draw[-{Stealth[length=2.6mm]}, thick] (knot) -- (jp)
        node[midway, above, font=\scriptsize] {Kauffman / Burau};

  \draw[-{Stealth[length=2.2mm]}, thick, gray] (cpx.north) to[bend left=18] (ph.south west);
\end{tikzpicture}

\caption{pytop hesaplama haritası: girdiden invaryanta giden yollar.}
\end{figure}

%-------------------------------------------------------------------
\section{Hesaplamalı Çekirdek: İlk Bakış}

\subsection*{Örnek 9.1 — Sonlu Bir Çember Modeli ve $\pi_1 = \mathbb{Z}$}

\begin{lstlisting}[language=Python]
fc = finite_circle()
r = pi1_space(fc)
print("uzay adi   :", r.space_name)
print("grup tipi  :", r.group_type)
print("uretec say :", len(r.presentation.generators))
print("baginti say:", len(r.presentation.relators))
print("serbest mi :", r.is_free())
print("H1 betti   :", r.abelianization_betti)
\end{lstlisting}

\begin{verbatim}
uzay adi   : S^1
grup tipi  : infinite_cyclic
uretec say : 1
baginti say: 0
serbest mi : True
H1 betti   : 1
\end{verbatim}

Yalnız 4 nokta, hiç bağıntısı olmayan tek üreteç: $\pi_1 = \langle a \rangle
= \mathbb{Z}$. Sonlu bir uzay, sürekli bir çemberle aynı temel grubu taşır.

\subsection*{Örnek 9.2 — Üçgen Kenarından Çemberin Homolojisi}

\begin{lstlisting}[language=Python]
circle = simplicial_complex([[0], [1], [2], [0, 1], [1, 2], [0, 2]])
for k in (0, 1):
    h = simplicial_homology(circle, k)
    print(f"H{k}: betti={h.betti} torsion={h.torsion}")
print("betti_numbers:", betti_numbers(circle))
\end{lstlisting}

\begin{verbatim}
H0: betti=1 torsion=()
H1: betti=1 torsion=()
betti_numbers: (1, 1)
\end{verbatim}

$H_0 = \mathbb{Z}$ (tek bileşen), $H_1 = \mathbb{Z}$ (bir delik) --- çemberin
klasik Betti imzası $(1, 1)$, ham simpleks listesinden SNF yoluyla hesaplandı.

\begin{dikkat}
\texttt{simplicial\_complex} yüz-kapalı girdi bekler: kenarları
\texttt{[0,1]} verip köşeleri \texttt{[0]} eklemeyi unutursanız hata alırsınız.
Üçgeni \texttt{[0,1,2]} olarak eklerseniz içi dolar ve $H_1 = 0$ olur.
\end{dikkat}

\subsection*{Örnek 9.3 — Nokta Bulutundan Kalıcı Homoloji}

\begin{lstlisting}[language=Python]
import math

pts = tuple(
    (round(math.cos(2 * math.pi * k / 8), 4),
     round(math.sin(2 * math.pi * k / 8), 4))
    for k in range(8)
)
space = FiniteMetricSpace(carrier=pts, distance=math.dist)
pairs = persistent_homology(space, max_dimension=1, max_scale=1.0)

essential = [p for p in pairs if p.is_essential]
print("toplam cift :", len(pairs))
print("kalici H0   :", sum(1 for p in essential if p.dimension == 0))
print("kalici H1   :", sum(1 for p in essential if p.dimension == 1))
print("betti       :", persistence_betti_numbers(pairs))
\end{lstlisting}

\begin{verbatim}
toplam cift : 9
kalici H0   : 1
kalici H1   : 1
betti       : {0: 1, 1: 1}
\end{verbatim}

Kalıcı sınıflar yine $(1, 1)$'i verir: bir bileşen ve bir delik. Sadece sıralı
noktalardan, hiçbir bağlantı bilgisi vermeden, pytop çemberin şeklini geri
kurtarır.

%-------------------------------------------------------------------
\section{Alıştırmalar}

\subsection*{Kodlama}

\begin{enumerate}[label=K\arabic*.]
\item \texttt{simplicial\_complex} ile içi \textbf{dolu} üçgeni
      (\texttt{[0,1,2]} facet'i, tüm yüzleriyle) kurun ve \texttt{betti\_numbers}
      çıktısını boş üçgeninkiyle karşılaştırın. $H_1$ neden kayboldu?
\item Örnek 9.3'teki nokta sayısını 8'den 4'e düşürün. Kalıcı $H_1$ hâlâ
      $1$ mi? Sonucu çemberin örnekleme yoğunluğuyla ilişkilendirin.
\end{enumerate}

\subsection*{Teori}

\begin{enumerate}[label=T\arabic*.]
\item Sonlu bir uzayın (Örnek 9.1) sürekli bir çemberle aynı $\pi_1$'e sahip
      olması nasıl mümkün? McCord teoreminin ne söylediğini bir cümleyle açıklayın.
\end{enumerate}
